\documentclass[letterpaper]{article}
\usepackage[submission]{aaai2027}

% Keep the supplementary PDF in the same official AAAI format as the paper.
\usepackage[hyphens]{url}
\usepackage{graphicx}
\urlstyle{rm}
\def\UrlFont{\rm}
\usepackage{natbib}
\usepackage{caption}
\frenchspacing
\usepackage{amsmath,amssymb}
\usepackage{array}
\usepackage{booktabs}

\graphicspath{
  {../../../CoFiTok-internal/artifacts/figures/ablation_hyperparameters_2026-07-11/}
}

\pdfinfo{
/TemplateVersion (2027.1)
}

\setcounter{secnumdepth}{0}
\renewcommand{\thetable}{S\arabic{table}}
\renewcommand{\thefigure}{S\arabic{figure}}

\newcommand{\method}{CoFiTok}
\newcommand{\epshat}{\hat{\epsilon}}
\newcommand{\xzero}{x_0}
\newcommand{\xt}{x_t}
\newcommand{\zk}{z^{(k)}}
\newcommand{\sk}{S_k}

\title{Supplementary Material for \method: Coarse-to-Fine Denoising Tokens for Pixel-Space Diffusion}
\author{Anonymous Submission}
\affiliations{}

\begin{document}
\maketitle
\raggedbottom

\section{Supplementary Scope}

This document separates comparison evidence from the main paper's
mechanism-centered tables by protocol, not by outcome. P0 contains
matched-dataset and matched-step pixel-space training. P1 contains
official-pretrained tokenizer reconstruction. A third tier contains official
ImageNet-256 50K evaluations for architecture-incompatible methods. Values from
different tiers are not used to claim direct superiority.

\begin{table*}[t]
  \centering
  \scriptsize
  \caption{Comparison protocols used throughout the paper and supplement.}
  \label{tab:supp-protocol-map}
  \resizebox{\textwidth}{!}{%
  \begin{tabular}{llllp{0.34\textwidth}}
    \toprule
    Tier & Methods & Coverage & Evidence & Valid interpretation \\
    \midrule
    P0 matched training & \method, direct dense, endpoint-only, ablations, Improved DDPM, EDM & Eight datasets & Train/eval under locked per-dataset steps & Within-row short-budget context; external methods are not compute matched \\
    P1 reconstruction & FlexTok, TiTok & Eight datasets & Official-pretrained reconstruction on 256 images & Tokenizer reconstruction context only; not a generation comparison \\
    Official related & D-AR, MAR, ReTok & ImageNet-256 & Official-pretrained 50K evaluation & Secondary large-model context; not matched retraining against \method \\
    \bottomrule
  \end{tabular}%
  }
\end{table*}

\section{Complete Evidence Coverage}

Table~\ref{tab:supp-coverage-matrix} exposes the full eight-dataset by
15-method Cartesian product. A blocked cell marks a cross-task comparison that
is not interchangeable with the available official evaluation; official
results are never counted as matched training.

\begin{table*}[t]
  \centering
  \tiny
  \caption{Evidence coverage matrix (8 datasets $\times$ 15 methods). T = completed matched train/eval; E = completed official-pretrained reconstruction eval-only; B = matched cross-task protocol blocked; O = separate official ImageNet-256 50K eval-only result. B+O retains the blocked matched cell while exposing the secondary official result.}
  \label{tab:supp-coverage-matrix}
  \resizebox{\textwidth}{!}{%
  \begin{tabular}{lcccccccc}
    \toprule
    Method & CIFAR & Tiny-IN & IN64-HF & IN64-strict & FFHQ & AFHQv2 & IN256-10\% & IN256 \\
    \midrule
    \method & T & T & T & T & T & T & T & T \\
    Direct dense epsilon & T & T & T & T & T & T & T & T \\
    Endpoint-only factorized & T & T & T & T & T & T & T & T \\
    Legacy clean-prefix objective & T & T & T & T & T & T & T & T \\
    No denoise-path prefix term & T & T & T & T & T & T & T & T \\
    Add clean-space monotonic term & T & T & T & T & T & T & T & T \\
    Simultaneous prediction & T & T & T & T & T & T & T & T \\
    Deep-$S_k$ synthesis & T & T & T & T & T & T & T & T \\
    Improved DDPM & T & T & T & T & T & T & T & T \\
    EDM & T & T & T & T & T & T & T & T \\
    D-AR & B & B & B & B & B & B & B & B+O \\
    FlexTok & E & E & E & E & E & E & E & E \\
    TiTok & E & E & E & E & E & E & E & E \\
    MAR & B & B & B & B & B & B & B & B+O \\
    ReTok & B & B & B & B & B & B & B & B+O \\
    \bottomrule
  \end{tabular}%
  }
\end{table*}

The locked matrix contains 80 completed P0 train/eval cells, 16 P1 eval-only
cells, and 24 protocol-blocked cells. Three official 50K rows remain separate.

\section{External-Method Results}

The next four tables are generated from the locked evidence JSON. They report
P0 generation, FlexTok/TiTok reconstruction, official D-AR/MAR/ReTok
ImageNet-256 results, and the paper's claim boundary.

\newpage

\section{How to Read the External Comparisons}

\paragraph{Improved DDPM and EDM.}
Table~\ref{tab:p0-generation-current} is the direct external pixel-diffusion
context. Dataset, resolution, split, and nominal optimizer steps are aligned,
but parameter count, batch size, nominal images seen, and NFE remain explicit.
EDM is the low-resolution Frechet-style winner on seven of eight rows, and the
direct dense predictor wins the remaining row. These results rule out a broad
unconditional-generation superiority claim for \method. They do not test the
paper's distinctive interface claim because monolithic predictors do not expose
ordered component prefixes.

\paragraph{FlexTok and TiTok.}
Table~\ref{tab:p1-tokenizer-recon-current} uses official-pretrained models and
evaluates reconstruction at $256\times256$ on 256 images per dataset. These
methods tokenize image content; \method{} factorizes a dense diffusion noise
prediction. Their reconstruction PSNR and Frechet-style values therefore
document adjacent capability rather than a shared generation objective.

\input{../../../CoFiTok-internal/artifacts/reports/paper_evidence_report_2026-07-11_final/latex_tables/paper_tables_appendix.tex}

\paragraph{D-AR, MAR, and ReTok.}
Table~\ref{tab:official-related-eval-only} reports 50K-sample official
evaluations on ImageNet-256. The rows use available pretrained checkpoints and
their corresponding generation protocols. They are useful large-model context,
but they cannot be relabeled as eight-dataset matched-step retraining. The
Inception Score for MAR follows the ADM-style evaluator used in the recorded
run and is not mixed with the differently scaled value reported by another
evaluation implementation.

\section{Matched P0 Protocol Details}

The P0 protocol fixes dataset identity, spatial resolution, split, and nominal
optimizer steps. It does not pretend that different published architectures
have equal FLOPs. Table~\ref{tab:supp-p0-settings} exposes the remaining major
differences instead of hiding them in a single rank.

\begin{table*}[t]
  \centering
  \scriptsize
  \caption{Locked P0 short-budget settings. C = \method; D = direct dense; E = endpoint-only factorized. NFE is listed in the order C/D/Improved-DDPM/EDM.}
  \label{tab:supp-p0-settings}
  \resizebox{\textwidth}{!}{%
  \begin{tabular}{llrrrrrrl}
    \toprule
    Protocol family & Datasets & Resolution & Steps & C/D/E batch & DDPM/EDM batch & Quality images & Sample/real images & NFE \\
    \midrule
    CIFAR-10 & CIFAR-10 & 32 & 3000 & 64 & 32 & 512 & 1024/4096 & 50/50/50/79 \\
    64px & Tiny-IN, IN64-HF, IN64-strict, FFHQ, AFHQv2 & 64 & 5000 & 32 & 32 & 512 & 1024/4096 & 50/50/50/79 \\
    ImageNet-256 & 10\% and full train & 256 & 5000 & 4 & 4 & 256 & 512/2048 & 50/50/50/79 \\
    \bottomrule
  \end{tabular}%
  }
\end{table*}

\begin{table*}[t]
  \centering
  \scriptsize
  \caption{Trainable parameter counts in the P0 protocols. The direct dense predictor is the parameter-matched factorization control; Improved DDPM and EDM are external architecture baselines.}
  \label{tab:supp-parameter-counts}
  \begin{tabular}{lrrrrr}
    \toprule
    Resolution & \method & Direct dense & Dense delta vs. \method & Improved DDPM & EDM \\
    \midrule
    $32\times32$ & 81,808 & 80,707 & $-1.35\%$ & 9,593,539 & 3,870,531 \\
    $64\times64$ & 81,808 & 80,707 & $-1.35\%$ & 21,715,907 & 3,870,531 \\
    $256\times256$ & 15,776 & 15,690 & $-0.55\%$ & 7,200,579 & 952,291 \\
    \bottomrule
  \end{tabular}
\end{table*}

Only the direct dense control isolates factorization at nearly equal parameter
count. Improved DDPM and EDM test whether the short-budget generator is
competitive with established monolithic objectives, while endpoint-only
factorized prediction tests whether ordered prefix supervision contributes
beyond retaining $K$ token components.

\section{Dataset and Source Conditions}

The two ImageNet-64 aliases are intentionally distinct. Results labeled
ImageNet-64 HF use the HF-derived source; results labeled ImageNet-64 strict use
the exact downsampled ImageNet source. No result is silently renamed between
them.

\begin{table*}[t]
  \centering
  \scriptsize
  \caption{Dataset roles in the locked evidence matrix.}
  \label{tab:supp-datasets}
  \resizebox{\textwidth}{!}{%
  \begin{tabular}{lllp{0.37\textwidth}}
    \toprule
    Dataset label & Resolution & Split role & Source condition \\
    \midrule
    CIFAR-10 & 32 & Train/test & Canonical CIFAR-10 staging \\
    Tiny-IN & 64 & Train/validation & Tiny ImageNet-200 staging \\
    ImageNet-64 HF & 64 & Train/validation & HF-derived ImageNet-1K 64x64 source \\
    ImageNet-64 strict & 64 & Train/validation & Exact downsampled ImageNet source from Academic Torrents \\
    FFHQ-64 & 64 & Train/evaluation & Non-classification face-domain extension \\
    AFHQv2-64 & 64 & Train/evaluation & Non-classification animal-domain extension \\
    ImageNet-256 10\% & 256 & 10\% train/full validation & Deterministic subset derived from the staged 256x256 source \\
    ImageNet-256 & 256 & Full train/validation & Full staged 256x256 ImageNet source \\
    \bottomrule
  \end{tabular}%
  }
\end{table*}

\section{Metric Definitions}

For prefix budget $m$, the cumulative prediction and pixel-space estimate are
\begin{equation}
  \epshat^{(1:m)}=\sum_{k=1}^{m}S_k(z^{(k)}), \qquad
  \hat{x}_0^{(m)}=\frac{x_t-\sigma_t\epshat^{(1:m)}}{\alpha_t}.
\end{equation}
Let $e_m$ be the mean-squared error between $\hat{x}_0^{(m)}$ and the
preregistered denoise-path target at prefix $m$. The reported path AUC is the
equal-spaced trapezoid average
\begin{equation}
  \operatorname{AUC}_{\mathrm{path}}
  =\frac{1}{K-1}\sum_{m=1}^{K-1}\frac{e_m+e_{m+1}}{2}.
\end{equation}
Lower values indicate that intermediate prefixes follow the intended partial
denoising path. Endpoint MSE instead uses only $m=K$ and does not measure the
quality of intermediate prefixes.

For component energy proportions $p_k$, effective token count is
\begin{equation}
  K_{\mathrm{eff}}=\exp\!\left(-\sum_k p_k\log p_k\right).
\end{equation}
The zero-token ratio divides the energy produced by $S_k(0)$ by the energy of
the corresponding observed components. A bias-free restricted synthesis
operator satisfies this diagnostic exactly up to numerical precision.

\section{Repeated and Confirmatory Evidence}

The matched 20K repeat uses K8 models, seeds 103 and 139, and 1,024 validation
images at $t=500$ on Tiny ImageNet-200 and ImageNet-64 HF. \method{} lowers path
AUC in all four paired dataset-seed comparisons, with a mean relative reduction
of 96.18\%. The mean endpoint $x_0$-MSE change is $+4.20\%$, so this repeat
supports prefix organization rather than an endpoint-quality win.

The ImageNet-256 K4 20K confirmatory run uses the same two seeds and exhaustively
evaluates all 24 component permutations. The learned order ranks first for both
seeds, the paired-bootstrap lower bounds for non-identity minus ordered AUC are
positive, and every permutation preserves the endpoint component sum up to
$3\times10^{-15}$ MSE. Relative to the parameter-matched direct dense control,
\method{} lowers endpoint MSE by 6.73\% on average in this confirmatory setting.

\section{Matched Component Ablations}

The component ablation suite uses ImageNet-64 HF, K8, 5,000 optimizer steps,
seeds 103/139, and 1,024 validation images at fixed $t=500$. All variants share
the data, model width, optimizer, and evaluation noise for each seed. The
progress-power sweep is scored against one common $p=1.5$ path; it is not scored
against a different target for each point.

\begin{table*}[t]
  \centering
  \small
  \caption{Matched component ablations (mean $\pm$ sample standard deviation).
  Lower is better for endpoint MSE, path AUC, and zero ratio; higher effective K
  indicates broader component use.}
  \label{tab:supp-component-ablation}
  \begin{tabular}{lrrrr}
    \toprule
    Variant & Endpoint $x_0$ MSE $\downarrow$ & Path AUC $\downarrow$ & Eff. K $\uparrow$ & Zero ratio $\downarrow$ \\
    \midrule
    \input{../../../CoFiTok-internal/artifacts/reports/aaai27_ablation_2026-07-11/component_ablation_rows.tex}\\
    \bottomrule
  \end{tabular}
\end{table*}

\begin{table*}[t]
  \centering
  \small
  \caption{Parameter-matched architecture control under the same two-seed 5k
  ImageNet-64 HF protocol. Direct dense has no native multi-prefix path.}
  \label{tab:supp-direct-dense-control}
  \begin{tabular}{lrrrr}
    \toprule
    Variant & Endpoint $x_0$ MSE $\downarrow$ & Path AUC $\downarrow$ & Eff. K $\uparrow$ & Zero ratio $\downarrow$ \\
    \midrule
    \input{../../../CoFiTok-internal/artifacts/reports/aaai27_ablation_2026-07-11/architecture_control_rows.tex}\\
    \bottomrule
  \end{tabular}
\end{table*}

The direct removals distinguish the path-prefix and path-component terms from
the endpoint objective. The simultaneous variant removes token-to-token
feedback while retaining the same heads and losses. The deep-synthesis variant
retains token-only inputs but introduces biased nonlinear layers, so its
zero-token ratio directly tests whether decoder capacity can carry a
token-independent prior. Full \method{} lowers path AUC by 93.0\% relative to
endpoint-only and by 45.7\%/37.9\% relative to removing the prefix/component
path term, respectively. Endpoint MSE changes by less than 0.2\% in the two
single-term comparisons. Sequential feedback changes the mean more modestly
($-2.8\%$ path AUC and $-1.0\%$ endpoint MSE) and wins path AUC on only one of
two seeds, so it is treated as an optional predictor implementation rather than
part of the core claim. Deep synthesis reaches lower path AUC but produces a
nonzero zero-token ratio of $0.0485\pm0.0050$.
The parameter-matched direct dense control is only 0.6\% lower in endpoint MSE
at 64px but exposes no multi-prefix interface; conversely, the ImageNet-256
confirmatory setting gives \method{} a 6.73\% endpoint advantage over direct
dense. Factorization is therefore supported as an interface with a small or
favorable endpoint tradeoff, not as a universal endpoint-quality guarantee.

\section{Hyperparameter Sensitivity}

Figure~\ref{fig:supp-ablation-hyperparameters} and
Table~\ref{tab:supp-hyperparameter-values} report the full two-seed sweeps. The
dashed lines mark the main values: $\lambda_p=0.15$, $\lambda_c=0.30$,
$p=1.5$, and $K=8$.
Increasing either loss weight improves path AUC with little endpoint movement;
the chosen weights are a moderate endpoint/path tradeoff. Power 1.5 is the
minimum under the common evaluation target. K8 has the lowest mean path AUC,
K4 the lowest endpoint MSE, and K16 no short-budget path advantage.

\begin{figure*}[t]
  \centering
  \includegraphics[width=\textwidth]{aaai27_ablation_hyperparameters.pdf}
  \caption{Matched ImageNet-64 HF hyperparameter sensitivity. Top row:
  denoise-path AUC; bottom row: endpoint $x_0$ MSE. Error bars are sample
  standard deviations over seeds 103/139. Lower is better.}
  \label{fig:supp-ablation-hyperparameters}
\end{figure*}

\begin{table*}[t]
  \centering
  \small
  \caption{Numerical values underlying
  Figure~\ref{fig:supp-ablation-hyperparameters}. Each entry is mean $\pm$
  sample standard deviation over seeds 103/139.}
  \label{tab:supp-hyperparameter-values}
  \begin{tabular}{lrrrr}
    \toprule
    Axis & Value & Endpoint $x_0$ MSE $\downarrow$ & Path AUC $\downarrow$ & Eff. K $\uparrow$ \\
    \midrule
    \input{../../../CoFiTok-internal/artifacts/reports/aaai27_ablation_2026-07-11/hyperparameter_sweep_rows.tex}\\
    \bottomrule
  \end{tabular}
\end{table*}

\begin{table*}[t]
  \centering
  \scriptsize
  \caption{Exploratory ImageNet-256 confirmatory generation controls. These rows do not replace the path and endpoint preregistered tests. Lower is better.}
  \label{tab:supp-imagenet256-generation-controls}
  \begin{tabular}{rlrrr}
    \toprule
    Seed & Method & Samples & Lowres Frechet & Inception Frechet \\
    \midrule
    103 & Endpoint-only factorized & 4096 & 11.1619 & 370.3672 \\
    103 & Direct dense & 4096 & 9.7396 & 449.3913 \\
    103 & \method & 4096 & 10.8232 & 349.4076 \\
    139 & Endpoint-only factorized & 4096 & 8.9753 & 404.6116 \\
    139 & Direct dense & 4096 & 10.5597 & 384.8126 \\
    139 & \method & 4096 & 11.0839 & 356.9311 \\
    \bottomrule
  \end{tabular}
\end{table*}

\section{Evidence Lock and Reproducibility}

All numerical tables imported above are generated from the final locked JSON
evidence package. The paper source does not hand-edit those metrics. The locked
state contains the per-run configuration, dataset alias, split, optimizer
steps, batch size, seed, parameter count, NFE, evaluation image count, and
sample count needed to distinguish the three protocol tiers. New experiments
must produce a new evidence package rather than silently overwrite these rows.

\end{document}
